\section{The Weierstrass form}
\subsection{Introduction}
The Weierstrass form is a canonical form of two matrices $A,B$ acting on the same space.
It is often referred to as generalized Jordan form. The motivation for Weierstrass form goes further than just generalizing Jordan and finding an appropriate basis - it plays crucial role in understanding linear systems of the form $$B\dot{x}=Ax$$
The relevance of such systems will be discussed in Chapter 5.

Let $A,B\in \lmap(V).$

Recalling the introduction of this paper, canonical forms are obtained by decomposing the space into invariant subspaces on which the operators have particularly simple matrix representations.

Assume that we are able to find $k$ subspaces $S_1,S_2,...,S_k$ such that each $S_i$ is $T$ and $L$-invariant, and $$\bigoplus_{i=1}^k S_k=V.$$Picking a basis $\beta_i$ for $S_i$ and composing $\bigcup\beta_i$ gives us a block-diagonal matrix representation of $A$ \textit{and} $B$, respectively, \textit{under the same basis}.

The problem is: such subspaces rarely exist. The maps $A$ and $B$ might be too independent of each other and therefore a subspace might act nice under one map, but arbitrarily under the other. For example, take
\[
    A=
    \begin{pmatrix}
        1  & 0         \\
        0  & 2  \\


    \end{pmatrix}
    ,
    B=
    \begin{pmatrix}
        0  & 1        \\
        1  & 0  \\


	\end{pmatrix}
\]
The line $W=\text{span}(e_1)$ is $A-$invariant: $w\in W \implies Aw=w.$

But, $w\in W\implies Bw\in \text{span}(e_2)$, so $W$ is not $B-$invariant.

In fact, for a pair $A,B$ there might be no non-trivial commonly invariant subspace.

Therefore, instead of attempting to describe the maps $A$ and $B$ via a specific basis, we do so with two different bases $\beta$ and $\gamma$ such that the matrices
\begin{equation}\label{form}
    _\beta[A]_ \gamma,\ _\beta[B]_ \gamma
\end{equation}
are of particularly nice type.
Prioritizing simplification, Weierstrass canonical form pursues this goal by having to put a restriction on $A$ and $B$.

We will define a property of the pairs $(A,B)$ referred to as \textit{regular}.

\begin{definition}
    A pair of maps $(A,B)$ is called \textit{regular} iff
    \begin{equation}\label{pencil}
        \exists \lambda\in \mathbb{C}\ s.t\ |A+\lambda B|\neq0
    \end{equation}
\end{definition}


\subsection{Statement}

We claim that:

For a regular pair of maps $(A,B)$, there exist bases $\beta,\gamma$ such that
\begin{align}\label{weierstrass}
    {}_{\beta}[A]_{\gamma}
    &=
    \begin{bmatrix}
        J & 0 \\
        0 & I \\
    \end{bmatrix}\\
    {}_{\beta}[B]_{\gamma}
    &=
    \begin{bmatrix}
        I & 0 \\
        0 & N \\
    \end{bmatrix}.
\end{align}

For the case where \(A\) and \(B\) are matrices, the statement is equivalent to:

There exists invertible \(P\) and \(Q\) s.t.

\begin{equation}
    \left(PAQ, PBQ\right) =
    \left(
        \begin{bmatrix}
            J & 0 \\
            0 & I \\
        \end{bmatrix},
        \begin{bmatrix}
            I & 0 \\
            0 & N \\
        \end{bmatrix},
    \right)
\end{equation}

Where \(I\) is the indentity matrix, \(J\) is a matrix in Jordan form and \(N\) is a nil-potent Jordan matrix.


\subsection{Matrix pencils}
A useful tool for studying the pairs \((A, B)\) where \(A, B \in \lmap(\vsV) \) is the so-called matrix pencil \footnote{For a rigorous analysis matrix pencils see \cite{gantmacher2026theory} page Chapter 12, page 24.} induced by the pair:

\begin{definition}
    The function \[P(x) = \penAB,  \ x \in \mathbb{C} \]
    is called the \textit{linear matrix pencil of the pair $(A,B)$}
\end{definition}


We say that a matrix pencil of the pair $(A,B)$ is \textit{regular} if and only if $(A,B)$ is regular.  It is called \textit{singular} otherwise.

The condition (\ref{pencil}) enables us to select a non-singular element of the matrix pencil which will help us derive a canonical form of (\ref{form}).

Before continuing, we will first derive some properties regarding the regularity of a matrix pencil.

\begin{lemma}
    \label{lemma:singular_from_vector}
    A matrix pencil $\penAB$ where $A, B \in \lmap(\vsV)$ is singular if and only if for all $x\in\C$, exists nonzero $\vv\in\vsV$ such that $(\penAB)\vv=\0$.
\end{lemma}
\begin{proof}
    \begin{align*}
        |\pen{A}{B}|=0
        &\iff \forall x\in\C,\, A+Bx=0\\
        &\iff \forall x\in\C,\, A+Bx \text{ is singular}\\
        &\iff \forall x\in\C,\, \exists \0\neq\vv\in\vsV,\, (A+Bx)\vv = \0\\
    \end{align*}
\end{proof}

\begin{lemma}
    \label{lemma:regular_symm}
    $A+Bx$ is regular if and only if $B+Ax$ is regular.
\end{lemma}
\begin{proof}
    Note that the forward direction is equivalent to the backward direction.

    If $A+Bx$ is regular, then there exists $x\in\C$ s.t. $|A+Bx|\neq0$.

    If $x=0$, then $A$ is invertible.
    \begin{equation*}
        |Ax+B|=|A(x\id+A\inv B)|=|A||x\id+A\inv B|
    \end{equation*}
    Pick $x\notin\spec -A\inv B$, we have $|Ax+B|\neq0$.

    Otherwise,
    \begin{equation*}
        |Ax\inv+B|=|x\inv(A+Bx)|=x^{-n}|A+Bx|\neq0
    \end{equation*}
\end{proof}

\begin{lemma}
    \label{lemma:regular_converge_inj}
    Let $A+Bx$ be a regular matrix pencil on $\vsV$, if $B\inv(A(\vsS))=\vsS$ for some subspace $\vsS\subseteq\vsV$, then $\ker_{\vsS}A=\{\0\}$.
\end{lemma}
\begin{proof}
    Suppose $\ker_{\vsS}A\neq\{\0\}$.

    Since $B\inv(A(\vsS))=\vsS$,
    \begin{equation*}
        \im_{\vsS}B = B(\vsS) = B(B\inv(A(\vsS))) \subseteq A(\vsS)
    \end{equation*}

    Hence, for any $x\in\C$, we have
    \begin{equation*}
        \im_{\vsS}(A+Bx)\subseteq\im_{\vsS}A+\im_{\vsS}Bx=\im_{\vsS}A
    \end{equation*}
    And
    \begin{equation*}
        \dim\ker_{\vsS}(A+Bx)=\dim\vsS-\dim\im_{\vsS}(A+Bx)\ge\dim\vsS-\dim\im_{\vsS}A=\dim\ker_{\vsS}A>0
    \end{equation*}

    Therefore $A+Bx$ is singular for all $x\in\C$.
    This contradicts the assumption for the matrix pencil to be regular.

    By contradiction, $\ker_{\vsS}A=\{\0\}$
\end{proof}


\subsection{Method } %TODO


Let $\vsV$ be a $n$-dimensional vector spaces.

Let $\pen{A}{B}$ be a regular matrix pencil, where $A, B$ are linear maps from $\vsV$ to $\vsV$.

% TODO: explain our goal with choosing subspaces
% TODO: Explain why in such form

Define
\begin{align} \label{eq:S_def}
    \vsS_0 &= \{\0\} \\
    \vsS_{m+1} &= B\inv(A(\vsS_{m})) \\
    \vsS &= \bigcup_{i\in\N} \vsS_i \\
\end{align}


\begin{align} \label{eq:W_def}
    \vsW_0 &= \vsV \\
    \vsW_{m+1} &= A\inv(B(\vsW_{m})) \\
    \vsW &= \bigcap_{i\in\N} \vsW_i \\
\end{align}

Note that here, $A$, $B$ and $A\inv$, $B\inv$ denote the image and pre-image, respectively.

% TODO: Observation that invertible has effects?

\begin{observation}\label{obv:wong_mono}
    For any \(m\in\N\), we have \(\vsS_m\subseteq\vsS_{m+1}\) and \(\vsW_m\supseteq\vsW_{m+1}\)
\end{observation}

Due to the finite-dimensionality of \(\vsV\), we conclude the following.
\begin{corollary}\label{col:wong_conv}
    There exists \(\kappa, \kappa'\in\N\) such that for all \(k\ge\kappa, k'\ge\kappa'\), we have
    \begin{align*}
        \vsS &= \vsS_k \\
        \vsW &= \vsW_{k'} \\
    \end{align*}
\end{corollary}

As a consequence of this, we have the following lemma:

\begin{lemma} \label{lemma:A_B_inj}
    \(A\) is injective on \(\vsS\) and \(B\) is injective on \(\vsW\).
\end{lemma}
\begin{proof}
    From Corollary \ref{col:wong_conv}, we have
    \begin{equation*}
        \vsS=\vsS_{\kappa+1}=B\inv(A(\vsS_{\kappa}))=B\inv(A(\vsS))
    \end{equation*}
    With this we may directly apply Lemma \ref{lemma:regular_converge_inj}, and conclude that \(A\) is injective on \(\vsS\).

    For the other case, recall that \((A,B)\) is regular if and only if \((B,A)\) is regular (Lemma \ref{lemma:regular_symm}).
    Therefore, we similarly conclude that \(B\) is injective on \(W\)

\end{proof}



Recall the definitions for $\vsS$ and $\vsW$.  $\vsS_{m+1} = B\inv(A(\vsS_{m}))$,
\begin{align*}
    \vv\in\vsS_{m+1}&\iff\exists\vu\in\vsS_{m}, B\vv=A\vu
\end{align*}
Hence, an alternative definition for $\vsS_m$ is
\begin{equation}\label{eq:S_def_alt}
    \vsS_m = \left\{\vv_m\in\vsV:\exists \vv_0,\vv_1,\dots,\vv_{m-1}\in\vsV,\, \vv_0 = \0 \land \forall i \in [m],\, B\vv_{i+1}=A\vv{i}\right\}
\end{equation}

Similarly, we have the alternative defination for \(\vsW_m\) as
\begin{equation}\label{eq:W_def_alt}
    \vsW_m = \left\{\vv_0\in\vsV:\exists \vv_1,\vv_2,\dots,\vv_m\in\vsV,\, \forall i \in [m],\, B\vv_{i+1}=A\vv_i\right\}
\end{equation}

We have redefined $\vsS$ and $\vsW$ sequentially due to having a more descrete view. Note that due to the injectiveness of \(A\) on \(\vsS\), we can make the following observation:
\begin{observation}
    For all $\vv\in\vsS$, there exists unique $\vu\in\vsS$ s.t. $B\vv=A\vu$.

    Hence, if for some vectors $\vv\in\vsS_{m+1}$ and $\vu\in\vsS$, $B\vv=A\vu$, then we have $\vu\in\vsS_m$
    % and if $\forall\vv\in\vsS_0,\vu\in\vsS$ s.t. $B\vv=A\vu$, we have $\vu\in\vsS_0=\{0\}$.
\end{observation}

Since $A\inv(B(\vsW))=\vsW$, $B$ is injective on $\vsW$.
So $\forall\vv\in\vsW$, $\exists!\vu\in\vsW$ s.t. $A\vv=B\vu$.

\subsection{Characterising \(\vsS\) and \(\vsW\) }


To show that \(\vsS\) and \(\vsW\) are a good division of \(\vsV\), we need to show \(\vsS\oplus\vsW=\vsV\).


To do so, we might want to take advantage of Lemma \ref{lem:img_sum_ker} by writing \(\vsS\) and \(\vsW\) as the extended kernel and image of some linear map \(X\).


In fact, Equation \ref{eq:S_def} is explicitly the definition of the extended kernel of the map  \(A\inv B,\) given that \(A\inv\) exists. Similarly,  Equation \ref{eq:W_def} is the generalized image of \(A\inv B\).
The only thing that stands in our way is that \(A\) might not be invertible.

For this purpose, we will introduce an equivalent definition for \(\vsS_m\) and \(\vsW_m\).

\begin{lemma}
    \label{lemma:alt_def_2}
    For any $\lambda \in \C$, we have
    \begin{align*}
        \vsS_{m+1} &= B\inv((A+\lambda B)(\vsS_m))\\
        \vsW_{m+1} &= (A+\lambda B)\inv(B(\vsW_m))\\
    \end{align*}
    % \begin{align*}
    %     \vsS&=\ker ((A+\lambda B)^{-1}B)^N\\
    %     \vsW&=\im ((A+\lambda B)^{-1}B)^N
    % \end{align*}
    % for some integer \(N\).
\end{lemma}
\begin{proof}

    Recall that
    \begin{align*}
        \vsS_m &= \{\vv_m:\exists \vv_0, \vv_1, \dots, \vv_{m-1},\,\vv_0=\0\land\forall i\in[m],\,B\vv_{i+1}=A\vv_i\}\\
        \vsW_m &= \{\vv_0:\exists \vv_1, \dots, \vv_m,\,\forall i\in[m],\,B\vv_{i+1}=A\vv_i\}\\
    \end{align*}
    Given sequence $\vv_0,\vv_1,\dots,\vv_m$, for any $\lambda\in\C$, one can construct $\vu_0,\vu_1,\dots,\vu_m$ defined as:
    \begin{equation*}
        \vu_i = \sum_{j=0}^{i}\begin{pmatrix}i\\j\end{pmatrix}\lambda^{i-j}\vv_j
    \end{equation*}
    We have
    \begin{align*}
        B\vu_{i+1}
        &=B\sum_{j=0}^{i+1}\begin{pmatrix}i+1\\j\end{pmatrix}\lambda^{i+1-j}\vv_j\\
        &=B\begin{pmatrix}i+1\\0\end{pmatrix}\lambda^{i+1}\vv_0+B\sum_{j=1}^{i}\begin{pmatrix}i+1\\j\end{pmatrix}\lambda^{i+1-j}\vv_j+B\begin{pmatrix}i+1\\i+1\end{pmatrix}\lambda^{0}\vv_{i+1}\\
        &=\lambda B\begin{pmatrix}i\\0\end{pmatrix}\lambda^{i}\vv_0+B\sum_{j=1}^{i}\left(\begin{pmatrix}i\\j\end{pmatrix}+\begin{pmatrix}i\\j-1\end{pmatrix}\right)\lambda^{i+1-j}\vv_j+\begin{pmatrix}i\\i\end{pmatrix}\lambda^{0}A\vv_{i}\\
        &=\lambda B\begin{pmatrix}i\\0\end{pmatrix}\lambda^{i+1}\vv_0+B\sum_{j=1}^{i}\begin{pmatrix}i\\j\end{pmatrix}\lambda^{i+1-j}\vv_j+B\sum_{j=1}^{i}\begin{pmatrix}i\\j-1\end{pmatrix}\lambda^{i+1-j}\vv_j+\begin{pmatrix}i\\i\end{pmatrix}\lambda^{0}A\vv_{i}\\
        &=\lambda B\begin{pmatrix}i\\0\end{pmatrix}\lambda^{i+1}\vv_0+B\sum_{j=1}^{i}\begin{pmatrix}i\\j\end{pmatrix}\lambda^{i+1-j}\vv_j+\sum_{j=1}^{i}\begin{pmatrix}i\\j-1\end{pmatrix}\lambda^{i+1-j}A\vv_{j-1}+\begin{pmatrix}i\\i\end{pmatrix}\lambda^{0}A\vv_{i}\\
        &=\lambda B\begin{pmatrix}i\\0\end{pmatrix}\lambda^{i+1}\vv_0+B\sum_{j=1}^{i}\begin{pmatrix}i\\j\end{pmatrix}\lambda^{i+1-j}\vv_j+\sum_{j=0}^{i-1}\begin{pmatrix}i\\j\end{pmatrix}\lambda^{i-j}A\vv_{j}+\begin{pmatrix}i\\i\end{pmatrix}\lambda^{0}A\vv_{i}\\
        &=\lambda B\sum_{j=0}^{i}\begin{pmatrix}i\\j\end{pmatrix}\lambda^{i-j}\vv_j+\sum_{j=0}^{i}\begin{pmatrix}i\\j\end{pmatrix}\lambda^{i-j}A\vv_{j}\\
        &=(A+\lambda B)\sum_{j=0}^{i}\begin{pmatrix}i\\j\end{pmatrix}\lambda^{i-j}\vv_j\\
        &=(A+\lambda B)\vu_i\\
    \end{align*}
    This process can also be reversed using $\lambda'=-\lambda$.
    Hence, there's a one-to-one relation between the sequence $\vv_i$ and $\vu_i$.
    %explain this more clearly

    Hence,
    \begin{align*}
        \vsS_m &= \{\vv_m:\exists \vv_0, \vv_1, \dots, \vv_{m-1},\,\vv_0=\0\land\forall i\in[m],\,B\vv_{i+1}=(A+\lambda B)\vv_i\}\\
        \vsW_m &= \{\vv_0:\exists \vv_1, \dots, \vv_m,\,\forall i\in[m],\,B\vv_{i+1}=(A+\lambda B)\vv_i\}\\
    \end{align*}
    or equivalently,
    \begin{align*}
        \vsS_{m+1} &= B\inv((A+\lambda B)(\vsS_m))\\
        \vsW_{m+1} &= (A+\lambda B)\inv(B(\vsS_m))\\
    \end{align*}
\end{proof}

With that in mind, due to the regularity of the pair \((A, B)\) We can choose some \(\lambda\in\C\) s.t. \(A+\lambda B\) is invertible and therefore we will consequently provide an explicit matrix \(X\) such that \[\vsS=\ker X,\ \vsW=im X\]
\begin{lemma} \label{lemma:S_W_sum}
    \(\vsS\oplus\vsW=\vsV\)
\end{lemma}
\begin{proof}
    Note that since \(\penAB\) is regular, there exist \(\lambda\in\C\) such that \(A+\lambda B\) is non-singular.

    Then, we can show
    \begin{align*}
        \vsS_m&=\ker((A+\lambda B)\inv B)^m\\
        \vsW_m&=\im((A+\lambda B)\inv B)^m\\
    \end{align*}
    by induction.

    For the base case, note that
    \begin{equation*}
        \vsS_0=\{\0\}=\ker((A+\lambda B)\inv B)^0
    \end{equation*}
    \begin{equation*}
        \vsW_0=\vsV=\im((A+\lambda B)\inv B)^0
    \end{equation*}

    Suppose this is true for some \(m\in\N\).
    We have
    \begin{align*}
        \vsS_{m+1}&=B\inv((A+\lambda B)(\vsS_m))\\
        &=\{\vv\in\vsV:Bv\in (A+\lambda B)(\vsS_m)\}\\
        &=\{\vv\in\vsV:(A+\lambda B)\inv Bv\in \vsS_m\}\\
        &=\{\vv\in\vsV:((A+\lambda B)\inv B)^{m+1}v=\0\}\\
        &=\ker((A+\lambda B)\inv B)^{m+1}\\
    \end{align*}
    and
    \begin{align*}
        \vsW_{m+1}&=((A+\lambda B)\inv(B(\vsS_m)))\\
        &=((A+\lambda B)\inv B)(\vsS_m)\\
        &=((A+\lambda B)\inv B)^{m+1}(\vsV)\\
        &=\im((A+\lambda B)\inv B)^{m+1}\\
    \end{align*}

    Therefore, we \ref{lem:img_sum_ker}, we have
    \begin{equation*}
        \vsV = \vsS\oplus\vsW
    \end{equation*}

\end{proof}


Having the decomposition $\vsV=\vsS \oplus \vsW$, we may construct a basis for $\vsV$ by uniting bases for $\vsS$ and $\vsW$,respectively.

The final form given at (\ref{weierstrass}) indicates that we should be searching for an additional basis for $\vsV$. The following lemma indicates how we might do that.



\begin{lemma}
    $A(\vsS) \oplus B(\vsW) = \vsV$.
\end{lemma}y
\begin{proof}
    First, we will show $A(\vsS)$ and $B(\vsW)$ are disjoint:

    Suppose $\vv\in A(\vsS)\cap B(\vsW)$.

    Then, there exists $\vu\in\vsS$, $\vu'\in\vsW$ s.t. $\vv=A\vu$, $\vv=B\vu'$.

    Hence, $A\vu=B\vu'$, $\vu\in A\inv(B\vu')\subseteq\vsW$.

    Since $\vsS$ and $\vsW$ are disjoint, $\vu\in\vsS\cap\vsW=\{\0\}$.
    \begin{equation*}
        \vv=A\vu=A\0=\0
    \end{equation*}

    This implies $A(\vsS)\cap B(\vsW) = \{\0\}$

    Now, note that \(A\) is injective on \(\vsS\) and \(B\) is injective on \(\vsW\).
    Hence, we have
    \begin{equation*}
        \dim(A(\vsS)\oplus B(\vsW))=\dim A(\vsS)+\dim B(\vsW)=\dim\vsS+\dim\vsW=n
    \end{equation*}
    So \(A(\vsS)+B(\vsW)=\vsV\)
\end{proof}

\subsection{Bases for the subspaces}\label{sectionbasis}

Now that we have characterised the subspaces, we are ready to choose appropriate bases. Before we begin, note that the maps \(A\) and \(B\) are well defined even on the constrained subspaces.
\begin{equation*}
    B(\vsS)=B(\vsS_{\kappa+1})=B(B\inv(A(\vsS_{\kappa})))\subseteq A(\vsS_\kappa)=A(\vsS)
\end{equation*}
and
\begin{equation*}
    A(\vsW)=A(\vsW_{\kappa'+1})=A(A\inv(B(\vsW_{\kappa'})))\subseteq B(\vsW_{\kappa'})=B(\vsW)
\end{equation*}

Hence, we can divide the domain into \(\vsS\) and \(\vsW\), and codomain into \(A(\vsS)\) and \(B(\vsW)\).
For each of those divisions, we can choose a suitable basis for \(A\) and \(B\)

\subsubsection{Basis for \(\vsS\) and \(A(\vsS)\)}\label{basisS}
First, we will only look at the domain of \(\vsS\) and the codomain of \(A(\vsS)\).

Now, we will try to choose a basis for \(\vsS\) and \(A(\vsS)\).

Note that after choosing a basis \(\beta_\vsS\) for \(\vsS\), the obvious choice of basis for \(A(\vsS)\) is \(\gamma_\vsS=A(\beta_\vsS)\).
This way, the matrix representation of \(A\) on \(\vsS\) will be the identity matrix \(I\).

We would like to also have a nice matrix representation for \(B\).

Note that since \(A\) is invertible and \({}_{\beta_\vsS}[A\inv]_{\gamma_\vsS}=I\),
\begin{equation*}
    {}_{\gamma_\vsS}[B]_{\beta_\vsS}={}_{\beta_\vsS}[A\inv]_{\gamma_\vsS}[B]_{\beta_\vsS}={}_{\beta_\vsS}[A\inv B]_{\beta_\vsS}
\end{equation*}

Hence, we can choose \(\beta_\{\vsS\}\) to be the basis where \(A\inv B\) is the simplest.
The obvious choice is the same method as we did jordan canonical form.

We claim \(A\inv B\) is nilpotent.
\begin{proof}
    Let $\vv_k\in\vsS=\vsS_k$.

    Define $\vv_i=A\inv B\vv_{i+1}$ for all $i\in[k]$.
    Note that
    \begin{align*}
        \vv_i&=A\inv B\vv_{i+1}\\
        A\vv_i&=B\vv_{i+1}\\
    \end{align*}
    So $\forall i\in[k]$, $\vv_i\in\vsS_i$.
    Notably, $\vv_0\in\vsS_0=\{\0\}$.

    Hence, $(A\inv B)^k\vv_k=\vv_0=\0$.

    Since this is true for any $\vv_k\in\vsS$, $(A\inv B)^k=0$.
    Hence, $A\inv B$ is nilpotent.
\end{proof}

Now, we can apply similar method as we did for Jordan canonical form.
For the subspace $\vsS$ and map $A\inv B$, we can find a basis composed of null-chains.
Let $$\{\vv_{1,1},\vv_{1,2},\dots,\vv_{1,m_1}\}, \{\vv_{2,1},\vv_{2,2},\dots,\vv_{2,m_2}\}, \dots, \{\vv_{t,1},\vv_{t,2},\dots,\vv_{t,m_{t}}\}$$ be such a basis.
% Let $\vv_0=\0$ and $\{\vv_1,\vv_2,\dots,\vv_m\}$ be such a chain.
We have $$\vv_{i,j}=A\inv B\vv_{i,j+1}$$.

Let $\vu_{i,j}=A\vv_{i,j}$.
\(\vu\) forms a basis for \(A(\vsS)\).

We have
\begin{equation*}
    B\vv_{i,j+1}=AA\inv B\vv_{i,j+1}=A\vv_{i,j}=\vu_{i,j}
\end{equation*}

Hence, \(B\) can be written in the form

 \begin{equation}\label{finish1}
    {}_{\gamma_{\vsS}}[B]_{\beta_\vsS}=
    \begin{bNiceMatrix}
        \begin{array}{c c c c}
            0 & 1 &        &   \\
              & 0 & \ddots &   \\
              &   & \ddots & 1 \\
              &   &        & 0 \\
        \end{array}
         &&&\text{\Huge 0} \\
        &
        \begin{array}{c c c c}
            0 & 1 &        &   \\
              & 0 & \ddots &   \\
              &   & \ddots & 1 \\
              &   &        & 0 \\
        \end{array}
         && \\
        &&\ddots&\\
        \text{\Huge 0}&&&
        \begin{array}{c c c c}
            0 & 1 &        &   \\
              & 0 & \ddots &   \\
              &   & \ddots & 1 \\
              &   &        & 0 \\
        \end{array} \\
    \CodeAfter
        \OverBrace[shorten,yshift=3pt]{1-1}{1-1}{m_1}
        \OverBrace[shorten,yshift=3pt]{2-2}{2-2}{m_2}
        \OverBrace[shorten,yshift=3pt]{4-4}{4-4}{m_{t}}
    \end{bNiceMatrix}
\end{equation}



\subsubsection{Basis for \(\vsW\) and \(B(\vsW)\)}\label{basisW}
Now we'll look at \(A\) and \(B\) as function from \(\vsW\) to \(B(\vsW)\).

We can apply the same method as above. However, \(B\inv A\) is no longer nilpotent.

However, we can use the conclusion from the proof for jordan canonical form.

For each eigenvalue of \(B\inv A\), we can look at the generalized eigenspace of \(B\inv A\).
The direct sum of the eigenspaces is equal to \(\vsW\).

For each eigenvalue \(\lambda_i\), we can find a basis  $$\beta_{\vsW}=\{\vv_{1,1},\vv_{1,2},\dots,\vv_{1,m_1}\}, \{\vv_{2,1},\vv_{2,2},\dots,\vv_{2,m_2}\}, \dots, \{\vv_{t,1},\vv_{t,2},\dots,\vv_{t,m_{t}}\}$$ where
% Let $\vv_0=\0$ and $\{\vv_1,\vv_2,\dots,\vv_m\}$ be such a chain.
$$\vv_{i,j}=(B\inv A - \lambda_i\id) \vv_{i,j+1}$$.

Let $\vu_{i,j}=B\vv_{i,j}$.
\(\vu\) form a basis $\gamma_{\vsW}$ for \(B(\vsW)\).

We have
\begin{equation*}
    A\vv_{i,j+1}=BB\inv A\vv_{i,j+1}=B\vv_{i,j}+\lambda_i B\vv_{i,j+1}=\vu_{i,j}+\lambda_i \vu_{i,j+1}
\end{equation*}

\begin{equation}\label{finish2}
    {}_{\gamma_\vsW}[A]_{\beta_\vsW}=
    % \begin{array}{cc c}
    %     &1&2_2\\
    %     \multirow{2}{0em}{$\left[\right.$}&&1\\
    %     &&2\\
    % \end{array}
    \begin{bNiceMatrix}
        \begin{array}{c c c c}
            \delta_0 & 1 &        &   \\
              & \delta_0 & \ddots &   \\
              &   & \ddots & 1 \\
              &   &        & \delta_0 \\
        \end{array}
         &&&\text{\Huge 0} \\
        &
        \begin{array}{c c c c}
            \delta_1 & 1 &        &   \\
              & \delta_1 & \ddots &   \\
              &   & \ddots & 1 \\
              &   &        & \delta_1 \\
        \end{array}
         && \\
        &&\ddots&\\
        \text{\Huge 0}&&&
        \begin{array}{c c c c}
            \delta_l & 1 &        &   \\
              & \delta_l & \ddots &   \\
              &   & \ddots & 1 \\
              &   &        & \delta_l \\
        \end{array} \\
    \CodeAfter
        \OverBrace[shorten,yshift=3pt]{1-1}{1-1}{k_0}
        \OverBrace[shorten,yshift=3pt]{2-2}{2-2}{k_1}
        \OverBrace[shorten,yshift=3pt]{4-4}{4-4}{k_{l}}
    \end{bNiceMatrix}
\end{equation}

Therefore, choosing \(\beta=\beta_\vsS \cup \beta_\vsW\) as basis for domain and \(\gamma=\gamma_\vsS \cup \gamma_\vsW\) as basis for co-domain,
We have

\begin{equation}\label{weierstrass}
\begin{aligned}
{}_{\gamma}[A]_{\beta}
&=
\begin{pmatrix}
J & 0\\
0 & I
\end{pmatrix},
\\[1em]
{}_{\gamma}[B]_{\beta}
&=
\begin{pmatrix}
I & 0\\
0 & N
\end{pmatrix}.
\end{aligned}
\end{equation}

% Before we begin, note the following necessary lemma about how the images of $A$ and $B$ on each of $S$ and $W$ intertwine.
% \begin{lemma}
%     \begin{equation}\label{incl1}
%     B(S)\subseteq A(S)
% \end{equation}
% and
% \begin{equation}\label{incl2}
%     A(W)\subseteq B(W)
% \end{equation}
% \end{lemma}
% \begin{proof}
%    The result follows due to Lemma \ref{charA} and due to Lemma \ref{charB}.
% \end{proof}


% This is useful in the following way of describing $_{\gamma}[A]_{\beta}$ and $_{\gamma}[B]_{\beta}$.\\ \\
% Since $\text{span}(\gamma_S)=A(S)$ by definition of $\gamma_S$, (\ref{incl1}) implies each element of $B(\beta_S)$ may be represented only in terms of $\gamma_S$.

% Similarly, $\text{span}(\gamma_W)=B(W)$ by definition of $\gamma_W$ and from (\ref{incl2}) each element of $A(\beta_W)$ is expressed only in terms of $\gamma_W$

% As a result, we may write ${}_{\gamma}[A]_{\beta}$ and ${}_{\gamma}[B]_{\beta}$ in the following way
% \begin{equation}\label{begginingoftheend}
%     \begin{aligned}
% {}_{\gamma}[A]_{\beta}
% &=
% \begin{pmatrix}
% _{\gamma_W}[A]_{\beta_W} & 0\\
% 0 & _{\gamma_S}[A]_{\beta_S}
% \end{pmatrix},
% \\[1em]
% {}_{\gamma}[B]_{\beta}
% &=
% \begin{pmatrix}
% _{\gamma_W}[B]_{\beta_W} & 0\\
% 0 & _{\gamma_S}[B]_{\beta_S}
% \end{pmatrix}.
% \end{aligned}
% \end{equation}


% Recall the definitions of $\gamma_S$ and $\gamma_W$.
% \begin{equation*}
% \gamma_S=A(\beta_S)=\{Av_{0}^{(1)},Av_{0}^{(2)},\dots,Av_{0}^{(m_0)}\} \cup \{Av_{1}^{(1)},Av_{1}^{(2)},\dots,Av_{1}^{(m_1)}\} \cup \dots \cup \{Av_{t}^{(1)},Av_{t}^{(2)},\dots,Av_{t}^{(m_t)}\}
%  \end{equation*}
%  \begin{equation*}
% \gamma_W=B(\beta_W)=\{Bu_{0}^{(1)},Bu_{0}^{(2)},\dots,Bu_{0}^{(k_0)}\} \cup \{Bu_{1}^{(1)},Bu_{1}^{(2)},\dots,Bu_{1}^{(k_1)}\} \cup \dots \cup \{Bu_{l}^{(1)},Bu_{l}^{(2)},\dots,Bu_{l}^{(k_l)}\}
%  \end{equation*}

% Since $\gamma_S=A(\beta_S)$, we may trivially conclude the matrix representation
% \begin{equation}\label{!1}
% {}_{\gamma_S}[A]_{\beta_S}=I_{|\beta_S|}
% \end{equation}


% Similarly, $\gamma_W=B(\beta_W)$ and therefore
% \begin{equation}\label{!2}
% _{\gamma_W}[B]_{\beta_W}=I_{|\beta_W|}
% \end{equation}

% As a consequence of (\ref{begginingoftheend}), we are left to determine a closed form for $_{\gamma_S}[B]_{\beta_S}$ and $_{\gamma_W}[A]_{\beta_W}$.\\\\
% Start with $_{\gamma_S}[B]_{\beta_S}$. Recall the definition of $\beta_S$ at (\ref{BAdel2}):

% \begin{equation}
%     \beta_S=\{v_{0}^{(1)},v_{0}^{(2)},\dots,v_{0}^{(m_0)}\} \cup \{v_{1}^{(1)},v_{1}^{(2)},\dots,v_{1}^{(m_1)}\} \cup \dots \cup \{v_{t}^{(1)},v_{t}^{(2)},\dots,v_{t}^{(m_t)}\}
%  \end{equation}
% \begin{equation}\label{important}
%    Bv_{i}^{(1)}=0 \text{ and } Bv_{i}^{(k+1)}=Av_{i}^{(k)}
% \end{equation}
% Therefore we may express $B(\beta_S)$ as
% \begin{equation}
%     B(\beta_S)=\{Bv_{0}^{(1)},Bv_{0}^{(2)},\dots,Bv_{0}^{(m_0)}\} \cup \{Bv_{1}^{(1)},Bv_{1}^{(2)},\dots,Bv_{1}^{(m_1)}\} \cup \dots \cup \{Bv_{t}^{(1)},Bv_{t}^{(2)},\dots,Bv_{t}^{(m_t)}\}
%  \end{equation}
% which by (\ref{important}) can be expressed in terms of $A$ in the following way.
% \begin{equation}
%     B(\beta_S)=\{0,Av_{0}^{(1)},\dots,Av_{0}^{(m_0-1)}\} \cup \{0,Av_{1}^{(1)},\dots,Av_{1}^{(m_1-1)}\} \cup \dots \cup \{0,Av_{t}^{(1)},\dots,Av_{t}^{(m_t-1)}\}
%  \end{equation}

%  We know that $\gamma_S=A(\beta_S)$ and thus we have that $_{\gamma_S}[B]_{\beta_S}$=$N$ where

%  \begin{equation}\label{finish1}
%     N=
%     % \begin{array}{cc c}
%     %     &1&2_2\\
%     %     \multirow{2}{0em}{$\left[\right.$}&&1\\
%     %     &&2\\
%     % \end{array}
%     \begin{bNiceMatrix}
%         \begin{array}{c c c c}
%             0 & 1 &        &   \\
%               & 0 & \ddots &   \\
%               &   & \ddots & 1 \\
%               &   &        & 0 \\
%         \end{array}
%          &&&\text{\Huge 0} \\
%         &
%         \begin{array}{c c c c}
%             0 & 1 &        &   \\
%               & 0 & \ddots &   \\
%               &   & \ddots & 1 \\
%               &   &        & 0 \\
%         \end{array}
%          && \\
%         &&\ddots&\\
%         \text{\Huge 0}&&&
%         \begin{array}{c c c c}
%             0 & 1 &        &   \\
%               & 0 & \ddots &   \\
%               &   & \ddots & 1 \\
%               &   &        & 0 \\
%         \end{array} \\
%     \CodeAfter
%         \OverBrace[shorten,yshift=3pt]{1-1}{1-1}{m_0}
%         \OverBrace[shorten,yshift=3pt]{2-2}{2-2}{m_1}
%         \OverBrace[shorten,yshift=3pt]{4-4}{4-4}{m_{t}}
%     \end{bNiceMatrix}
% \end{equation}
% We now derive a closed form of $_{\gamma_W}[A]_{\beta_W}$.
% Recall the definition of $\beta_W$ at (\ref{definitionbw}).
% \begin{equation}
% \beta_W=\{u_{0}^{(1)},u_{0}^{(2)},\dots,u_{0}^{(k_0)}\} \cup \{u_{1}^{(1)},u_{1}^{(2)},\dots,u_{1}^{(k_1)}\} \cup \dots \cup \{u_{l}^{(1)},u_{l}^{(2)},\dots,u_{l}^{(k_l)}\}
%  \end{equation}
%   \begin{equation}\label{AinB}
%         Au_{i}^{(1)}=\delta_i Bu_{i}^{(1)}\text{ and }Au_{i}^{(k+1)}=\delta_i Bu_{i}^{(k+1)}+Bu_{i}^{(k)}
%     \end{equation}

% Expressing $A(\beta_W)$ gives
% \begin{equation}
% A(\beta_W)=\{Au_{0}^{(1)},Au_{0}^{(2)},\dots,Au_{0}^{(k_0)}\} \cup \{Au_{1}^{(1)},Au_{1}^{(2)},\dots,Au_{1}^{(k_1)}\} \cup \dots \cup \{Au_{l}^{(1)},Au_{l}^{(2)},\dots,Au_{l}^{(k_l)}\}
%  \end{equation}

% We can express each vector of $A(\beta_W)$ in terms of $B$, due to (\ref{AinB}), as
% \begin{align*}
% A(\beta_W)=\{\delta_0 Bu_{0}^{(1)},\ \delta_0 Bu_{0}^{(2)}+Bu_{0}^{(1)},\dots,\delta_0 Bu_{0}^{(k_0)}+Bu_{0}^{(k_0-1)}\} \cup\\
% \{\delta_1 Bu_{1}^{(1)},\ \delta_1Bu_{1}^{(2)}+Bu_{1}^{(1)},\dots,\delta_1 Bu_{1}^{(k_1)}+Bu_{1}^{(k_1-1)}\} \cup\\
% \vdots\\
% \{\delta_l Bu_{l}^{(1)},\ \delta_lBu_{l}^{(2)}+Bu_{l}^{(1)},\dots,\delta_l Bu_{l}^{(k_l)}+Bu_{l}^{(k_l-1)}\}\quad
%  \end{align*}

% Because $\gamma_W=B(\beta_W)$, we have that $_{\gamma_W}[A]_{\beta_W}=J$ where \\ \\
% \begin{equation}\label{finish2}
%     J=
%     % \begin{array}{cc c}
%     %     &1&2_2\\
%     %     \multirow{2}{0em}{$\left[\right.$}&&1\\
%     %     &&2\\
%     % \end{array}
%     \begin{bNiceMatrix}
%         \begin{array}{c c c c}
%             \delta_0 & 1 &        &   \\
%               & \delta_0 & \ddots &   \\
%               &   & \ddots & 1 \\
%               &   &        & \delta_0 \\
%         \end{array}
%          &&&\text{\Huge 0} \\
%         &
%         \begin{array}{c c c c}
%             \delta_1 & 1 &        &   \\
%               & \delta_1 & \ddots &   \\
%               &   & \ddots & 1 \\
%               &   &        & \delta_1 \\
%         \end{array}
%          && \\
%         &&\ddots&\\
%         \text{\Huge 0}&&&
%         \begin{array}{c c c c}
%             \delta_l & 1 &        &   \\
%               & \delta_l & \ddots &   \\
%               &   & \ddots & 1 \\
%               &   &        & \delta_l \\
%         \end{array} \\
%     \CodeAfter
%         \OverBrace[shorten,yshift=3pt]{1-1}{1-1}{k_0}
%         \OverBrace[shorten,yshift=3pt]{2-2}{2-2}{k_1}
%         \OverBrace[shorten,yshift=3pt]{4-4}{4-4}{k_{l}}
%     \end{bNiceMatrix}
% \end{equation}

% Therefore, (\ref{!1}),  (\ref{!2}),  (\ref{finish1}) and  (\ref{finish2}) together imply that
% \begin{equation*}
% \begin{aligned}
% {}_{\gamma}[A]_{\beta}
% &=
% \begin{pmatrix}
% J & 0\\
% 0 & I
% \end{pmatrix},
% \\[1em]
% {}_{\gamma}[B]_{\beta}
% &=
% \begin{pmatrix}
% I & 0\\
% 0 & N
% \end{pmatrix}.
% \end{aligned}
% \end{equation*}
% With this, we have given a closed Weierstrass form, as desired. The above form may be easily concluded for matrices, where we would have transformation invertible matrices $P$ and $Q$ which change the coordinates to $\gamma$ and $\beta$, respectively.





% Therefore, $\dim\vsS+\dim\vsW=n$, so $\dim\vsT+\dim\vsTp=n$.
% Since $\vsT$ and $\vsTp$ are disjoint, $\vsV=\vsT\oplus\vsTp$.

% Note that $A$ is injective from $\vsS$ to $\vsT$, and since $\vsT=A(\vsS)$, it is also surjective.
% So we can define $A\inv:\vsT\to\vsS$ s.t. $AA\inv\vv=\vv$ and $A\inv A\vu=\vu$ for all $\vv\in\vsT$ and $\vu\in\vsS$.
% Similarly, we can define $B\inv:\vsTp\to\vsW$.

% Note that
% \begin{equation*}
%     B(\vsS)=B(\vsS_{k+1})=B(B\inv(A(\vsS_{k})))\subseteq A(\vsS_k)=\vsT
% \end{equation*}
% and
% \begin{equation*}
%     A(\vsW)=A(\vsW_{k'+1})=A(A\inv(B(\vsW_{k'})))\subseteq B(\vsW_{k'})=\vsTp
% \end{equation*}

% So we can have $B':\vsS\to\vsS=A\inv\circ B$ and $A':\vsW\to\vsW=B\inv\circ A$, being well-defined on the domain.

% Note that $B'$ is nilpotent.
% \begin{proof}
%     Let $\vv_k\in\vsS=\vsS_k$.

%     Define $\vv_i=B'\vv_{i+1}$ for all $i\in[k]$.
%     Note that
%     \begin{align*}
%         \vv_i&=B'\vv_{i+1}\\
%         \vv_i&=A\inv B\vv_{i+1}\\
%         A\vv_i&=B\vv_{i+1}\\
%     \end{align*}
%     So $\forall i\in[k]$, $\vv_i\in\vsS_i$.
%     Notably, $\vv_0\in\vsS_0=\{\0\}$.

%     Hence, $(B')^k\vv_k=\vv_0=\0$.

%     Since this is true for any $\vv_k\in\vsS$, $(B')^k=0$.
%     Hence, $B'$ is nilpotent.
% \end{proof}

% We also claim that $\spec -A'=\spec(A, B)$. Which we will show using double inclusion lemma.

% \begin{proof}
%     If $x\in\spec -A'$, then we have $|x\id+A'|=0$.
%     Hence, $\exists\vv\in\vsW\subseteq\vsV$ s.t. $(x\id+A')\vv=\0$.
%     \begin{align*}
%         (x\id+A')\vv=\0
%         &\implies B(x\id+A')\vv=\0\\
%         &\implies (BA'+Bx)\vv=\0\\
%         &\implies (A+Bx)\vv=\0\\
%         &\implies |A+Bx|=0\\
%     \end{align*}

%     If $x\in\spec(A, B)$, then we have $|A+Bx|=0$.
%     Hence, there exists vector $\vv\in\vsV$ s.t. $(A+Bx)\vv=\0$.
%     And we have $A\vv=-xB\vv$

%     Note that
%     \begin{equation*}
%         \vsW_{m+1}=(A+xB)\inv(B(\vsS_m))\supseteq(A+xB)\inv(B(\{\0\}))=\ker(A+xB)\ni\vv
%     \end{equation*}
%     Hence,
%     \begin{align*}
%         (x\id+A')\vv
%         &=x\vv+B\inv A\vv\\
%         &=x\vv-xB\inv B\vv\\
%         &=\0\\
%     \end{align*}
% \end{proof}


% Now, we can apply similar method as we did for Jordan canonical form.
% For the subspace $\vsS$ and map $B'$, we can find a basis composed of null-chains.
% Let $$\{\vv_{0,1},\vv_{0,2},\dots,\vv_{0,m_0}\}, \{\vv_{1,1},\vv_{1,2},\dots,\vv_{1,m_1}\}, \dots, \{\vv_{\gamma-1,1},\vv_{\gamma-1,2},\dots,\vv_{\gamma-1,m_{\gamma-1}}\}$$ be such a basis.
% % Let $\vv_0=\0$ and $\{\vv_1,\vv_2,\dots,\vv_m\}$ be such a chain.
% We have $$\vv_{i,j}=B'\vv_{i,j+1}=A\inv B\vv_{i,j+1}$$.
% Let $\vu_{i,j}=A\vv_{i,j}$. Since $A$ is an invertible map from $\vsS$ to $\vsT$, $\vu_{i,j}$ form a basis for $\vsT$.
% We also have
% \begin{equation*}
%     B\vv_{i,j+1}=AA\inv B\vv_{i,j+1}=A\vv_{i,j}=\vu_{i,j}
% \end{equation*}

% Hence, using $\vv$ and $\vu$ as basis, we can write
% \begin{align*}
%     \leftindex_\vu[A]_\vv&=I\\
%     \leftindex_\vu[B]_\vv&=N\\
% \end{align*}
% where
% \vspace{2\baselineskip}
% % \setcounter
% \begin{equation*}
%     N=
%     % \begin{array}{cc c}
%     %     &1&2_2\\
%     %     \multirow{2}{0em}{$\left[\right.$}&&1\\
%     %     &&2\\
%     % \end{array}
%     \begin{bNiceMatrix}
%         \begin{array}{c c c c}
%             0 & 1 &        &   \\
%               & 0 & \ddots &   \\
%               &   & \ddots & 1 \\
%               &   &        & 0 \\
%         \end{array}
%          &&&\text{\Huge 0} \\
%         &
%         \begin{array}{c c c c}
%             0 & 1 &        &   \\
%               & 0 & \ddots &   \\
%               &   & \ddots & 1 \\
%               &   &        & 0 \\
%         \end{array}
%          && \\
%         &&\ddots&\\
%         \text{\Huge 0}&&&
%         \begin{array}{c c c c}
%             0 & 1 &        &   \\
%               & 0 & \ddots &   \\
%               &   & \ddots & 1 \\
%               &   &        & 0 \\
%         \end{array} \\
%     \CodeAfter
%         \OverBrace[shorten,yshift=3pt]{1-1}{1-1}{m_0}
%         \OverBrace[shorten,yshift=3pt]{2-2}{2-2}{m_1}
%         \OverBrace[shorten,yshift=3pt]{4-4}{4-4}{m_{\gamma-1}}
%     \end{bNiceMatrix}
% \end{equation*}

% Similarly, for the map $A'$ on $\vsW$, we can also find a simple basis.
% However, since $A'$ is not nilpotent, we need to apply the method used for Jordan form.

% We can write $\vsW$ as the direct sum of extended eigenspace for $A'$, and find a basis for each eigenspace.
% This way, for each eigen value $\lambda$, we can find chains where for each chain $\{\vv_1,\vv_2,\dots,\vv_m\}$, we have $\vv_i=(A'-\lambda\id)\vv_{i+1}$.

% Let $\vu_i=B\vv_i$. With similar argument as before, $\vu$ form a basis for the eigenspace.
% Note that
% \begin{equation*}
%     A\vv_{i+1}=BB\inv A\vv_{i+1}=BA'\vv_{i+1}=B(A'-\lambda\id)\vv_{i+1}+\lambda B\vv_{i+1}=B\vv_i+\lambda B\vv_{i+1}=\lambda\vu_{i+1} + \vu_i
% \end{equation*}

% Hence, using $\vv$ and $\vu$ as basis, we can write
% \begin{align*}
%     \leftindex_\vu[A]_\vv&=J\\
%     \leftindex_\vu[B]_\vv&=I\\
% \end{align*}

